\begingroup
\setlength{\abovedisplayskip}{6pt}
\setlength{\belowdisplayskip}{6pt}
\setlength{\abovedisplayshortskip}{4pt}
\setlength{\belowdisplayshortskip}{4pt}
\setlength{\intextsep}{6pt}

The calculations follow the course notes (pp. 132--135). Pressures are interpreted as root-mean-square (RMS) values and intensities as time averages for a progressive plane wave.

\textbf{1. Input data.} The sound speed and air density are
\[
\boxed{c=340\,\mathrm{m/s},\qquad \rho=1.2\,\mathrm{kg/m^3}.}
\]
The specified sound pressure, loudspeaker diameter, and sound levels are
\[
\boxed{p=55\,\mathrm{mPa}=0.055\,\mathrm{Pa},\qquad
D=1.55\,\mathrm{m},\qquad L_{\mathrm{sp}}=78\,\mathrm{dB},\qquad
L_{\mathrm{horn}}=118\,\mathrm{dB}.}
\]

\textbf{2. Acoustic impedance.} The plane-wave acoustic impedance is
\[
Z=\rho c=(1.2)(340)=408\,\mathrm{Pa\,s/m}
\]
so that
\[
\boxed{Z=408\,\mathrm{rayl}.}
\]

\textbf{3. Sound intensity for $p=55\,\mathrm{mPa}$.} Using the plane-wave pressure--intensity relation,
\[
I=\frac{p^2}{Z}=\frac{(0.055)^2}{408}
  =7.414\times10^{-6}\,\mathrm{W/m^2}.
\]
Therefore,
\[
\boxed{I=7.41\,\mu\mathrm{W/m^2}.}
\]

\textbf{4. Reference intensity and pressure.} The reference intensity is
\[
\boxed{I_0=10^{-12}\,\mathrm{W/m^2}=1\,\mathrm{pW/m^2}.}
\]
For the specified air properties, the consistent reference pressure is
\[
p_0=\sqrt{ZI_0}
    =\sqrt{(408)(10^{-12})}
    =2.0199\times10^{-5}\,\mathrm{Pa}.
\]
Thus,
\[
\boxed{p_0=20.199\,\mu\mathrm{Pa}\approx20\,\mu\mathrm{Pa}.}
\]

\textbf{5. Decibel values for the specified pressure.} The intensity level is
\[
L_I=10\log_{10}\left(\frac{I}{I_0}\right)
 =10\log_{10}\left(\frac{7.414\times10^{-6}}{10^{-12}}\right)
 \approx68.7007\,\mathrm{dB}.
\]
The pressure level is
\[
L_p=20\log_{10}\left(\frac{p}{p_0}\right)
 =20\log_{10}\left(\frac{0.055}{2.0199\times10^{-5}}\right)
 \approx68.7007\,\mathrm{dB}.
\]
Therefore,
\[
\boxed{L_I=L_p=68.7\,\mathrm{dB}.}
\]

\textbf{6. Equality of the two decibel values.} Since $I=p^2/Z$ and $I_0=p_0^2/Z$,
\[
\frac{I}{I_0}=\frac{p^2/Z}{p_0^2/Z}=\left(\frac{p}{p_0}\right)^2.
\]
Taking the base-10 logarithm gives
\[
10\log_{10}\left(\frac{I}{I_0}\right)
 =10\log_{10}\left[\left(\frac{p}{p_0}\right)^2\right]
 =20\log_{10}\left(\frac{p}{p_0}\right).
\]
The intensity and pressure formulas therefore describe the same physical sound level. The unrounded value $p_0=\sqrt{ZI_0}$ is used in the calculations so that the equality is exact for $Z=408\,\mathrm{rayl}$; $20\,\mu\mathrm{Pa}$ is the usual rounded nominal reference.

\textbf{7. Loudspeaker at $78\,\mathrm{dB}$.} The intensity corresponding to the loudspeaker level is
\[
I_{\mathrm{sp}}=I_0 10^{L_{\mathrm{sp}}/10}
 =10^{-12}10^{78/10}
 =6.3096\times10^{-5}\,\mathrm{W/m^2}.
\]
The corresponding RMS pressure is
\[
p_{\mathrm{sp}}=\sqrt{ZI_{\mathrm{sp}}}
 =\sqrt{(408)(6.3096\times10^{-5})}
 =0.1604\,\mathrm{Pa}.
\]
For a circular loudspeaker face,
\[
A=\frac{\pi D^2}{4}=\frac{\pi(1.55)^2}{4}=1.8869\,\mathrm{m^2}.
\]
Assuming the stated intensity is uniform and normal to this circular area, the sound power is
\[
P_{\mathrm{sp}}=I_{\mathrm{sp}}A
 =(6.3096\times10^{-5})(1.8869)
 =1.1906\times10^{-4}\,\mathrm{W}.
\]
The physical quantities are therefore
\[
\boxed{I_{\mathrm{sp}}=63.1\,\mu\mathrm{W/m^2},\qquad
p_{\mathrm{sp}}=160.4\,\mathrm{mPa},\qquad
P_{\mathrm{sp}}=119.1\,\mu\mathrm{W}.}
\]

\textbf{8. Pressure from a $118\,\mathrm{dB}$ car horn.} The horn intensity and pressure are
\[
I_{\mathrm{horn}}=I_0 10^{118/10}=0.631\,\mathrm{W/m^2},
\]
\[
p_{\mathrm{horn}}=p_0 10^{118/20}
 =20.199\times10^{-6}10^{118/20}
 =16.04\,\mathrm{Pa}.
\]
Thus,
\[
\boxed{p_{\mathrm{horn}}=16.04\,\mathrm{Pa}.}
\]

\textbf{9. Pain threshold comparison.} The course-note pain reference values are
\[
\boxed{I_{\mathrm{pain}}=1\,\mathrm{W/m^2},\qquad
p_{\mathrm{pain}}=\sqrt{ZI_{\mathrm{pain}}}\approx20.2\,\mathrm{Pa}.}
\]
The corresponding intensity level is
\[
L_{\mathrm{pain}}=10\log_{10}\left(\frac{1}{10^{-12}}\right)=120\,\mathrm{dB}.
\]
The $78\,\mathrm{dB}$ loudspeaker level is $42\,\mathrm{dB}$ below this reference, while the $118\,\mathrm{dB}$ horn level is $2\,\mathrm{dB}$ below it. Numerically, the horn pressure is close to the course-note pain-pressure reference, whereas the loudspeaker pressure is much smaller.

\textbf{10. Graduate work.} Sound intensity is proportional to the square of pressure,
\[
I=\frac{p^2}{Z}.
\]
Consequently, the intensity ratio is the square of the pressure ratio. The logarithm converts this square into a factor of two,
\[
10\log_{10}\left(\frac{I}{I_0}\right)
 =10\log_{10}\left[\left(\frac{p}{p_0}\right)^2\right]
 =20\log_{10}\left(\frac{p}{p_0}\right).
\]
This is why intensity uses $10\log_{10}$ while pressure uses $20\log_{10}$. Both formulas give the same decibel value when their reference quantities satisfy $I_0=p_0^2/Z$.

\clearpage
\matlabcaption{MATLAB calculations for B.6.}
\lstinputlisting[style=hwmatlab]{matlab/B6.m}
\endgroup
